\section{Introduction}
\label{sec:intro}
% ============================================================

Large language model (LLM)-based coding agents such as Claude Code~\cite{claudecode2025}, Codex~\cite{openaicodex2021}, and OpenCode~\cite{opencode2026} combine language models with tools, code execution, and operating-system access. Conventional code assistants such as GitHub Copilot~\cite{githubcopilot2021} suggest code for human approval. Coding agents can also inspect repositories, edit files, invoke compilers and tests, execute shell commands, and debug failures iteratively~\cite{yang2024sweagent,wang2025openhands}. They now operate in development environments that contain credentials, dependency managers, build systems, and version-control state~\cite{hou2024large}.

Coding agents perform these actions with the developer's privileges and also process untrusted project content~\cite{greshake2023indirect,zhan2024injecagent,ruan2024toolemu}. A malicious GitHub issue can therefore instruct an agent diagnosing a build failure to fetch and execute a remote script, and the agent may comply~\cite{guo2024redcode,andriushchenko2024agentharm}. Because agent-generated tool calls may be executed without human approval, a harmful action can reach the execution boundary before a human intervenes. Runtime monitors can still block such candidate actions, but doing so requires inspecting each one along the trajectory. We instead study an earlier control point: placing security policy in the context that generates every action, so the model can be steered before a harmful call is produced. The need is substantial. Across six LLMs, undefended agents comply with 46--80\% of malicious code-execution requests, and DeepSeek-V3.2~\cite{deepseekai2025deepseekv32} also produces malware with an average judge score of 7.76 out of 10~\cite{gu2024survey}.

Existing defenses address this risk at four stages of the agent pipeline (Figure~\ref{fig:overview}). At the \emph{model level}, alignment~\cite{ouyang2022training,bai2022training,rafailov2024dpo} requires access to model weights and costly retraining, which excludes API-only deployers. At the \emph{input stage}, safety filters~\cite{metallamaguard3,metapromptguard2025} can reject harmful user requests, but they cannot protect against malicious instructions encountered later during the agent's interaction with repository content or tool output. At the \emph{execution stage}, runtime verifiers~\cite{luo2025agrail,jia2025task} inspect candidate actions before execution and provide a stronger enforcement boundary, but they add a model or checking component to the tool-use loop. In \emph{prompt space}, defenses place security instructions in the agent context. Existing methods use boundary reminders, provenance delimiters, or authentication tags to distinguish untrusted data from instructions~\cite{yi2025benchmarking,hines2024defending,wang2024fath}. These methods use general rules across tasks and leave open how threat-specific policies should be synthesized and scoped within a finite prompt budget. We investigate this question for API-only deployments. The system prompt can govern action formation throughout the multi-step loop and be updated without retraining or an additional runtime component.

%We introduce \sysname, a prompt-space defense that synthesizes \emph{security skills} and places them in a coding agent's system prompt. A skill is a natural-language instruction document based on the emerging Agent Skills standard~\cite{agentskills2025}. \sysname can synthesize a skill from known attacks through \emph{proactive mining} or from recorded agent failures through \emph{reactive learning}. The resulting skill remains in the system prompt throughout the session. Operators also determine its coverage before deployment. A general-purpose agent uses one \emph{all-classes} skill. An environment exposed to a known family of mechanisms can use a \emph{per-bundle} skill. A narrowly constrained environment with one known threat class can use a \emph{per-class} skill. All configurations use the same prompt budget and require no request-time classification or routing.

We therefore propose \sysname, a prompt-space defense that injects \emph{security skills} into a coding agent's system prompt. A skill is a natural-language instruction document that extends agent behavior without modifying model weights~\cite{agentskills2025}. \sysname derives security skills from structured attack data. \emph{Proactive mining} extracts observed attack patterns, whereas \emph{reactive learning} captures the framings behind recorded agent failures. Each skill is injected in full at session start and remains active throughout the interaction. Because the prompt budget is fixed, increasing coverage forces more threat classes to share the same space, reducing the detail available to each class, while narrowing scope preserves class-specific detail at the cost of covering fewer threats. To capture this tradeoff, we study three deployment scopes that reflect how much threat information is known before a session begins. General-purpose agents use an \emph{all-classes} skill, deployments exposing a known mechanism family use a corresponding \emph{per-bundle} skill, and \emph{per-class} skills apply when the exact threat class is known in advance and provide an upper bound on fine-grained provisioning. All configurations use the same prompt budget and require no request-time classification or routing.

We evaluate \sysname on RedCode~\cite{guo2024redcode}, the only publicly available and established coding-agent benchmark that jointly measures harmful tool execution and malware generation. Four questions organize the evaluation. \textbf{Q1: Are prompt-space security skills effective?} Yes. The all-classes skill reduces execution attack success rate (ASR) from 67.4\% to 43.6\% and malware-generation severity from 3.37 to 0.58, while per-class provisioning further lowers execution ASR to 14.5\% when the threat class is known in advance. \textbf{Q2: What limits coverage in one finite prompt?} Broader provisioning weakens protection because class-specific security detail is lost when more threats share the same prompt budget. Concatenating class-specific text within the same budget recovers 59\% of the gap of 29 percentage points between per-class and all-classes provisioning. \textbf{Q3: Does a fixed policy remain effective under non-adaptive jailbreaks?} Partially. Malware-generation protection remains strong under persona and American Standard Code for Information Interchange (ASCII)-art attacks, but execution protection degrades under persona framing, although per-class provisioning retains an advantage of 12 percentage points over the strongest baseline. \textbf{Q4: Does prevention increase benign refusal?} Little. Across 731 SWE-Bench Pro task descriptions~\cite{deng2025swe}, safety-grounded refusal remains at or below 0.68\% in 32 of the 36 configurations we evaluated.

We provide the first systematic study of fixed-budget security policies synthesized from attack data for a coding agent's system prompt, using agent skills as the packaging format. Our work contributes security-policy synthesis and provisioning. Our code will be released publicly upon publication. The main contributions are listed below.
\begin{itemize}[leftmargin=*, itemsep=2pt]
\item \textbf{We develop \sysname, a prompt-space first line of defense for API-only coding agents.} \sysname synthesizes security policy offline from known attacks or recorded failures and places the policy in the operator-controlled system prompt. It requires no model-weight changes, auxiliary model, or runtime check.

\item \textbf{We identify deployment scope and synthesis quality as key determinants of prompt-space security.} Narrower scopes strengthen in-scope protection, while directly combining class-specific policy text recovers much of the loss at broader scopes.

\item \textbf{We evaluate security across models, agents, and jailbreaks.} Across six LLMs and three harnesses, per-bundle skills surpass the strongest clean-input baseline. Under jailbreaks, the defense retains its malware-generation advantage and loses its execution advantage.

\item \textbf{We measure the benign refusal cost of prompt-space prevention.} Proactive skills introduce little safety-grounded refusal across 731 benign SWE-Bench Pro task descriptions. End-to-end benign task success remains outside the scope of our evaluation.
\end{itemize}
